\section{结论}
\label{sec:conclusion}

本文介绍了FruitTreeScanner系统的设计与实现，该系统利用消费级LiDAR设备进行果实产量估算。系统通过多模态融合策略结合基于LiDAR的3D点云采集与CoreML图像检测。关键贡献包括：自适应DBSCAN聚类算法解决了LiDAR数据中距离相关密度变化的问题；多模态融合方法利用了2D和3D检测的互补能力；以及双路线产量估算策略结合了基于体积和冠层回归的方法。

模拟数据实验结果表明该方法的技术可行性，在中等遮挡条件下检测精度超过80\%。然而，这些结果需要通过真实果园数据的田间实验进行验证，然后才能得出关于实际性能的结论。

当前工作的局限性包括评估完全依赖模拟数据，无法捕捉真实果园环境的全部复杂性。遮挡校正模型使用了可能不适用于所有冠层结构的简化统计假设。系统还面临可变光照条件、背景植被复杂性和果实外观变化带来的挑战。

未来工作方向包括在不同果实品种和果园条件下进行全面的田间验证试验、通过在多视角捕获上训练的基于学习的方法改进遮挡校正模型、优化移动设备上的实时处理算法以及与农场管理信息系统集成以进行实际部署。

\section*{致谢}
本工作呈现系统设计和可行性分析。所有实验数据有待田间验证。

\section*{数据可用性}
系统实现代码和处理算法基于FruitTreeScanner iOS应用程序。本研究中使用的评估数据是模拟的，将在未来的田间研究中补充真实果园数据。